\SourcePage{354}
\SGAsection{6}{Chern 环}
本节的基环取为一个固定的\Term{binomial-ring}{二项式环} $K$（2.7）。我们使用\Term{category}{范畴} $\mathscr C$（1.1），因而回到第 1 节的记号：将 $\widehat G$ 留作定义在 $\mathscr C$ 上的\Term{functor}{函子}（1.1）的记号，并以 $\widehat G_{\circ i}$ 表示它在 $\mathscr C_0$ 上的\Term{restriction}{限制}（1.8）。

\textbf{6.1。} 在 3.5 和 3.6 中，我们定义了如下\Term{functor-morphism}{函子态射}：
\[
\begin{gathered}
(\widehat G_{\circ i})\times(\widehat G_{\circ i})
 \xrightarrow{\nabla}\widehat G_{\circ i},\\
\widehat G^t_{\circ i}\xrightarrow{\gamma'}
 (\widehat G^u_{\circ i})\circ(\widehat G^t_{\circ i})_{\nabla},\\
\widehat G_{\circ i}\xrightarrow{a}\widehat G_{\circ i},
\end{gathered}
\]
最后一个是关于运算法则 $\nabla$ 与 $\gamma_t(a)$ 相乘，其中 $a\in K$。

由 1.9，得到\Term{morphism}{态射}
\[
\begin{gathered}
\widehat G\times\widehat G\xrightarrow{\nabla}\widehat G_g,\\
\widehat G^t\xrightarrow{\gamma'}
 \widehat G^u_{\circ i}\circ\widehat G^t_{g\nabla},\\
\widehat G\xrightarrow{a}\widehat G_g.
\end{gathered}
\]
由\Term{polynomial}{多项式} $P'_n,Q'_{i,j},R_{n,a}$ 的\Term{isobaric-polynomial}{等权性}（3.5、3.6、2.9）可知，第一个和第三个\Term{functor}{函子}的像落在 $\widehat G$ 中，第二个的像则落在关于 $u$ 的\Term{formal-power-series}{形式幂级数}所成的\Term{subgroup}{子群}中，其中非零次项的系数都在 $\widehat G^t$ 中。

将第一个和第三个与\Term{canonical}{典范}\Term{morphism}{态射}
$\widehat G\longrightarrow1+\widehat{\operatorname{Gr}(\cdot)}^+$（1.11）复合，便由 1.12 得到 $\mathscr G$ 上的两个\Term{functor-morphism}{函子态射}：
\SourcePage{355}
\[
\begin{gathered}
(1+\widehat{\cdot}^{+})\times(1+\widehat{\cdot}^{+})
 \xrightarrow{*}1+\widehat{\cdot}^{+},\\
1+\widehat{\cdot}^{+}\xrightarrow{a}1+\widehat{\cdot}^{+}.
\end{gathered}
\]
它们在\Term{functor}{函子} $1+\widehat{\cdot}^{+}$ 上定义了\Term{multiplication}{乘法}和 $K$ 的运算。由 1.13，\Term{multiplication}{乘法}由它在
$(1+X,1+Y)\in(1+\widehat{K[X,Y]}^+)^2$ 上的值确定；将它记作 $*$，于是
\[
(1+X)*(1+Y)=\frac{1+X+Y}{(1+X)(1+Y)}.
\tag{6.1.1}
\]
这一运算法则的第 $n$ 个\Term{universal-polynomial}{泛多项式}是 $P'_n$ 的权 $n$ 的\Term{isobaric-polynomial}{等权分量}（1.15），记作 $P''_n$。

$K$ 的运算仍用指数记号表示（3.8）。定义这些运算的\Term{universal-polynomial}{泛多项式}是 $R_{n,a}$ 的最低次\Term{isobaric-polynomial}{等权分量}（1.15）；由于这些\Term{polynomial}{多项式}本身就是\Term{isobaric-polynomial}{等权的}，所得仍是 $R_{n,a}$ 自身。

设 $A\in\operatorname{Ob}(\mathscr G)$；由上述运算法则定义的环 $1+\widehat A^+$ \Term{generic}{一般}没有\Term{identity-element}{单位元}。因此，通过添加 $K$ 中的标量来扩充它，引入新环 $K\times(1+\widehat A^+)$，其运算法则如下：对任意 $a,a'\in A$ 及 $\varphi,\varphi'\in1+\widehat A^+$，
\[
\begin{aligned}
(a,\varphi)+(a',\varphi')&=(a+a',\varphi\varphi'),\\
(a,\varphi)\cdot(a',\varphi')
 &=(aa',\varphi^{a'}\cdot\varphi'^{a}\cdot(\varphi*\varphi')).
\end{aligned}
\tag{6.1.2}
\]
这样构造的环以 $(1,1)$ 为\Term{identity-element}{单位元}，而 $1+\widehat A^+$ 可等同于\Term{ideal}{理想} $(0)\times(1+\widehat A^+)$。它自然带有 $K$-\Term{algebraic-structure}{代数结构}，满足
\[
a\cdot(a',\varphi)=(a,1)\cdot(a',\varphi)
 =(aa',\varphi^a).
\tag{6.1.3}
\]
\SourcePage{356}
如下在它上面赋予 $K$-\Term{lambda-algebra}{$\lambda$-代数}结构。前面已经指出，\Term{functor}{函子}
$\gamma':\widehat G^t\longrightarrow\widehat G^u_{\circ i}\circ\widehat G^t_{g\nabla}$ 的像落在关于 $u$ 的\Term{formal-power-series}{形式幂级数}所成的\Term{subgroup}{子群}中，其非零次项的系数都在 $\widehat G^t$ 中。利用\Term{canonical}{典范}\Term{morphism}{态射}
$\widehat G\longrightarrow1+\widehat{\operatorname{Gr}(\cdot)}^+$，由它得到
$\widehat G\longrightarrow K\times(1+\widehat{\operatorname{Gr}(\cdot)}^+)$，最终得到\Term{morphism}{态射}
\[
\widehat G^t\xrightarrow{\gamma'}
\widehat G^u_{\circ i}\circ
 \bigl(K\times(1+\widehat{\operatorname{Gr}(\cdot)}^+)\bigr).
\]
根据 1.12，它定义了 $\mathscr G$ 上的一个\Term{functor-morphism}{函子态射}
\[
1+\widehat{\cdot}^{+}\longrightarrow
\widehat G^u_{\circ i}\circ\bigl(K\times(1+\widehat{\cdot}^{+})\bigr).
\]
定义这一\Term{morphism}{态射}的\Term{universal-polynomial}{泛多项式}是
$\gamma'^i(1+X_1+\cdots+X_n+\cdots)$ 的各个系数；第 $j$ 个系数是 $Q'_{i,j}$ 中权 $j$ 的\Term{isobaric-polynomial}{等权分量}（证明与 1.15 类似），记作 $Q''_{i,j}$。

最后，利用 $K$ 的 $\lambda$-结构，以自然的方式将这样构造的\Term{morphism}{态射}扩充到 $K\times(1+\widehat{\cdot}^{+})$。这样得到 $K\times(1+\widehat{\cdot}^{+})$ 上的\Term{pre-lambda-ring}{预 $\lambda$-环}结构，它由相应的 $\gamma$ 运算确定，满足
\[
\gamma_t(a,\varphi)=\gamma_t(a)\cdot
 \left(1+\sum_{p\geqslant1}\gamma'^p(\varphi)t^p\right).
\tag{6.1.4}
\]
由于不会混淆，以后将这些运算记作 $\gamma$，不用 $\gamma'$。

要证明这确实是\Term{lambda-ring}{$\lambda$-环}结构，只需验证（3.7.2）的\Term{commutativity}{交换性}。对于 $K$ 中的元素，这已经知道；对于 $1+\widehat{\cdot}^{+}$ 中的元素，则由 $\widehat G$ 的（3.7.2）的\Term{commutativity}{交换性}取商得到。所需结论随即成立。

\SourcePage{357}
\begin{SGAresult}{定义}{6.2}
6.1 中在\Term{graded-algebra}{分次 $K$-代数}的\Term{category}{范畴}上定义的、取值于 $K$-\Term{lambda-algebra}{$\lambda$-代数}\Term{category}{范畴}的\Term{functor}{函子}，称为\Term{chern-ring}{Chern 环}；它在 $A\in\operatorname{Ob}(\mathscr G)$ 上的值记作 $\operatorname{Ch}(A)$（可能混淆时记作 $\operatorname{Ch}_K(A)$）。
\end{SGAresult}
由公式（6.1.1）、（6.1.2）、（6.1.4）和（3.6.2），得到关系式
\[
\begin{gathered}
(1,1+X)\cdot(1,1+Y)=(1,1+X+Y),\\
\lambda^i(1,1+X)=0\qquad(i\geqslant2).
\end{gathered}
\tag{6.2.1}
\]
文献 [1] 正是用这些关系来定义\Term{chern-ring}{Chern 环}的。反过来，也可以由这些公式重新得到关系式（6.1.1）和（6.1.4）。

下面说明引入\Term{chern-ring}{Chern 环}的原因。置身于例 3.9 的情形，以 $A$ 表示一个交换\Term{local-scheme}{局部}\Term{ringed-topos}{环化拓扑斯}的 $K'$ 关于 3.10 中定义的\Term{filtration}{滤过}的\Term{associated-graded-ring}{相伴分次环}（为简便起见，假定这个 $K'$ 是\Term{lambda-ring}{$\lambda$-环}）；$K$ 是\Term{augmentation}{增广}的目标环。后面将会看到，对于\Term{augmented-lambda-ring}{增广 $\lambda$-环}的任意元素，可以定义取值于\Term{associated-graded-ring}{相伴分次环}中的\Term{chern-class}{Chern 类}（6.7）。因此，对每个\Term{locally-free-sheaf}{局部自由}\Term{finite-type}{有限型}\Term{sheaf}{层}，都能关联 $\operatorname{Ch}(A)$ 中的一个元素：其在因子 $K$ 上的\Term{set-component}{分量}是该\Term{sheaf}{层}的\Term{rank}{秩}，在因子 $1+\widehat A^+$ 上的\Term{set-component}{分量}则是\Term{chern-class}{Chern 类}之和。公式（6.2.1）于是就是通常的公式，给出两个\Term{invertible-sheaf}{可逆层}的\Term{tensor-product}{张量积}以及一个\Term{invertible-sheaf}{可逆层}的\Term{exterior-power}{外幂}的\Term{chern-class}{Chern 类}与\Term{rank}{秩}（也就是说，刚刚定义的从 $K'$ 到 $\operatorname{Ch}(A)$ 的映射是\Term{lambda-homomorphism}{$\lambda$-同态}）。

\SourcePage{358}
\textbf{6.3。} 对每个\Term{graded-algebra}{分次 $K$-代数} $A$，以 $A_{\mathbb Q}$ 表示 $K$-\Term{algebra}{代数} $A\otimes_{\mathbb Z}\mathbb Q$，赋予显然的\Term{grading}{分次}。根据 1.13，只需给出在 $1+X$ 上的值，就能定义一个具有\Term{functoriality}{函子性}的\Term{homomorphism}{同态}
$\operatorname{ch}:1+\widehat A^+\longrightarrow\widehat A_{\mathbb Q}^+$，取
\[
\operatorname{ch}(1+X)=\exp(X)-1.
\tag{6.3.1}
\]
这个\Term{homomorphism}{同态}与 6.1 中在 $1+\widehat A^+$ 上定义的\Term{multiplication}{乘法}相容：由 1.13，只需对 $1+X,1+Y\in1+\widehat{K[X,Y]}^+$ 证明这一点。而
\[
\begin{aligned}
\operatorname{ch}((1+X)*(1+Y))
 &=\operatorname{ch}\left(\frac{1+X+Y}{(1+X)(1+Y)}\right)\\
 &=\operatorname{ch}(1+X+Y)-\operatorname{ch}(1+X)-\operatorname{ch}(1+Y)\\
 &=\exp(X+Y)-\exp(X)-\exp(Y)+1\\
 &=\operatorname{ch}(1+X)\cdot\operatorname{ch}(1+Y).
\end{aligned}
\]
记 $\eta$ 为 $\widehat A_{\mathbb Q}^+$ 的\Term{endomorphism}{自同态}（关于 $A$ 具有\Term{functoriality}{函子性}），定义为
\[
\eta\left(\sum_{k\geqslant1}a_kt^k\right)
 =\sum_{k\geqslant1}\frac{(-1)^{k-1}}{(k-1)!}a_kt^k.
\]
那么，对任意 $\varphi\in1+\widehat A^+$，
\[
\operatorname{ch}(\varphi)=\eta(\operatorname{Log}\varphi).
\tag{6.3.2}
\]
事实上，$\operatorname{ch}$ 与 $\eta\circ\operatorname{Log}$ 是两个具有\Term{functoriality}{函子性}的\Term{abelian-group}{阿贝尔群}\Term{homomorphism}{同态}；只要它们在 $1+X$ 上一致，就相等，而这一点是显然的。

由此\Term{pushout}{推出} $\operatorname{ch}$ 与 $K$ 的运算交换。事实上，这里也只需对 $1+X$ 证明，而对任意 $a\in K$，有
\[
\begin{aligned}
\operatorname{ch}((1+X)^a)
 &=\operatorname{ch}(\exp(a\operatorname{Log}(1+X)))\\
 &=\eta(a\operatorname{Log}(1+X))\\
 &=a\eta(\operatorname{Log}(1+X))\\
 &=a\operatorname{ch}(1+X).
\end{aligned}
\]
\SourcePage{359}
因此，显然可以将 $\operatorname{ch}$ 扩充为 $K$-\Term{algebra}{代数}\Term{homomorphism}{同态}
\[
K\times(1+\widehat A^+)\longrightarrow K\oplus\widehat A_{\mathbb Q}^+.
\]
于是有
\[
\operatorname{ch}(a,\varphi)=a+\eta(\operatorname{Log}\varphi).
\tag{6.3.3}
\]
\begin{SGAresult}{定义}{6.4}
刚才定义的\Term{homomorphism}{同态}
$\operatorname{ch}:\operatorname{Ch}A\longrightarrow K\oplus\widehat A_{\mathbb Q}^+$ 称为\Term{chern-character}{Chern 特征}。
\end{SGAresult}
若 $A$ 的\Term{characteristic}{特征}为 $0$（即它是一个 $\mathbb Q$-\Term{algebra}{代数}），则\Term{chern-character}{Chern 特征}是\Term{isomorphism}{同构}。事实上，给定 $\psi\in K\oplus\widehat A^+$，可以写成 $\psi=a+\psi_1$，其中 $\psi_1$ 没有零次项；于是
$\psi\longmapsto(a,\exp(\eta^{-1}(\psi_1)))$ 定义了 $\operatorname{ch}$ 的逆\Term{morphism}{态射}。

\begin{SGAresult}{命题}{6.5}
设 $A$ 是\Term{graded-algebra}{分次 $K$-代数}；考虑 $\operatorname{Ch}(A)$ 上如下定义的\Term{filtration}{滤过}（\Term{degree-filtration}{次数滤过}）：
\[
(\operatorname{Ch}(A))_0=\operatorname{Ch}(A);\qquad
\forall n\geqslant1,\quad
(\operatorname{Ch}(A))_n=0\times(1+\widehat A^{\geqslant n}),
\tag{6.5.1}
\]
其中 $\widehat A^{\geqslant n}=\prod_{i\geqslant n}A^i$。有下列同余式。

\textup{i)} 设 $n\geqslant1$，$\varphi\in(\operatorname{Ch}A)_n$；置 $\varphi=1+a_n+\cdots$，则
\[
\operatorname{ch}(\varphi)=\frac{(-1)^{n-1}}{(n-1)!}a_n+\cdots.
\tag{6.5.2}
\]
\textup{ii)} 设 $\varphi\in(\operatorname{Ch}A)_m$，$\psi\in(\operatorname{Ch}A)_n$；置 $\varphi=1+a_m+\cdots$、$\psi=1+b_n+\cdots$，则
\[
\varphi*\psi
 =1-\frac{(m+n-1)!}{(m-1)!(n-1)!}a_mb_n+\cdots.
\tag{6.5.3}
\]
因而\Term{degree-filtration}{次数滤过}是\Term{ring-filtration}{环滤过}。
\end{SGAresult}
\begin{proof}
i）由构造，\Term{chern-character}{Chern 特征}是带\Term{filtration}{滤过}\Term{functor}{函子}的\Term{morphism}{态射}。由同余式
$\varphi\equiv1+a_n\pmod{(\operatorname{Ch}A)_{n+1}}$，
\SourcePage{360}
得到
$\operatorname{ch}(\varphi)\equiv\operatorname{ch}(1+a_n)\pmod{\widehat A^{\geqslant n+1}}$；而
\[
\operatorname{ch}(1+a_n)=\eta(\operatorname{Log}(1+a_n))
 =\frac{(-1)^{n-1}}{(n-1)!}a_n+\cdots.
\]
ii）通过代入，只需对 $A=K[X_i;Y_j]_{i,j\in\mathbb N}$ 及
$\varphi=1+\sum_{k\geqslant m}X_k$、$\psi=1+\sum_{k\geqslant n}Y_k$ 证明公式（6.5.3）。由于 $K$ \Term{torsion-free}{无挠}，可以将这个环\Term{embedding}{嵌入}它与 $\mathbb Q$ 的\Term{tensor-product}{张量积}中，从而看出 $\operatorname{ch}$ 是\Term{injectivity}{单射}。置 $\varphi*\psi=1+Z_p+\cdots$，其中 $Z_p$ 是 $\varphi*\psi$ 的第一个既非零又不等于 $1$ 的项。由 $\operatorname{ch}$ 的\Term{multiplication}{乘法}性及（6.5.2），有
\[
1+\frac{(-1)^{p-1}}{(p-1)!}Z_p+\cdots
 =\left(1+\frac{(-1)^{m-1}}{(m-1)!}X_m+\cdots\right)
  \left(1+\frac{(-1)^{n-1}}{(n-1)!}Y_n+\cdots\right),
\]
比较最低次项便得到 $Z_p$ 的值。
\end{proof}

\textbf{6.6。} \Term{lambda-ring}{$\lambda$-环} $\operatorname{Ch}A$ 向 $K$ \Term{augmentation}{增广}：事实上，$(a,\varphi)\longmapsto a$ 定义了从 $\operatorname{Ch}A$ 到 $K$ 的\Term{lambda-homomorphism}{$\lambda$-同态}。因此，它带有\Term{augmented-lambda-ring}{增广 $\lambda$-环}的\Term{canonical}{典范}\Term{filtration}{滤过}。下面将它与刚才定义的\Term{filtration}{滤过}作比较。

\begin{SGAresult}{命题}{6.6.1}
设 $x\in\operatorname{Ch}A$ 具有\Term{form}{形式} $x=(n,1+\cdots+a_n)$。则
\[
\begin{aligned}
\textup{i)}\quad&\lambda^i(x)=0\qquad(i\geqslant n);\\
\textup{ii)}\quad&\lambda^n(x)=(1,1+a_1);\\
\textup{iii)}\quad&\sum_{i=0}^n(-1)^i\lambda^i(x)
 =(0,1-(n-1)!a_n+\cdots),
\end{aligned}
\]
其中未写出的项的\Term{degree}{次数}都 $>n$。
\end{SGAresult}
\begin{proof}
由\Term{functoriality}{函子性}，只需对以下情形证明：
\SourcePage{361}
$A$ 是环 $K[X_1,\ldots,X_n]$，按 $X_i$ 的权赋予\Term{filtration}{滤过}，$X_i$ 的权为 $i$；所考虑的元素为 $(n,1+X_1+\cdots+X_n)$。

通过\Term{symmetric-function}{对称函数}将 $A$ \Term{embedding}{嵌入} $A'=K[T_1,\ldots,T_n]$，于是得到\Term{injectivity}{单射}\Term{homomorphism}{同态}
$\operatorname{Ch}A\longrightarrow\operatorname{Ch}A'$。有
\[
\begin{aligned}
(n,1+X_1+\cdots+X_n)
 &=\left(n,\prod_{i=1}^n(1+T_i)\right)\\
 &=\sum_{i=1}^n(1,1+T_i).
\end{aligned}
\]
由（6.2.1）\Term{pushout}{推出}
\[
\lambda_t(x)=\prod_{i=1}^n\bigl(1+(1,1+T_i)t\bigr),
\]
从而 $\lambda^k(x)=0$（$k>n$），并且
\[
\lambda^n(x)=\prod_{i=1}^n(1,1+T_i)
 =\left(1,1+\sum_{i=1}^nT_i\right)=(1,1+X_1).
\]
最后，
\[
\begin{aligned}
\sum_{i=1}^n(-1)^i\lambda^i(x)
 &=\lambda_{-1}(x)
 =\prod_{i=1}^n\bigl(1-(1,1+T_i)\bigr)\\
 &=(-1)^n\prod_{i=1}^n(0,1+T_i).
\end{aligned}
\]
根据（6.5.3），
\[
\prod_{i=1}^n(0,1+T_i)
 =\left(0,1+(-1)^{n-1}(n-1)!\prod_{i=1}^nT_i+\cdots\right),
\]
即
\[
\lambda_{-1}(x)=(0,1-(n-1)!X_n+\cdots).
\]
\end{proof}

\begin{SGAresult}{推论}{6.6.2}
关系式 6.6.1.iii）对每个\Term{augmentation}{增广}为 $n$ 的元素 $x$ 都成立。
\end{SGAresult}
为此，我们使用如下\Term{lemma}{引理}。

\begin{SGAresult}{\Term{lemma}{引理}}{6.6.3}
对每个 $n$，$(\operatorname{Ch}A)_n$ 是 $\operatorname{Ch}A$ 的\Term{lambda-ideal}{$\lambda$-理想}。
\end{SGAresult}
\begin{proof}
可以假定 $n\geqslant1$。设 $x=(0,\varphi)\in(\operatorname{Ch}A)_n$；可以写成
\[
\gamma^i(0,\varphi)=(0,1+Q''_{i,1}+\cdots+Q''_{i,j}+\cdots),
\]
其中 $Q''_{i,j}$ 是 $Q'_{i,j}$ 的权 $j$ 的\Term{isobaric-polynomial}{等权分量}（3.6）。因此，对每个 $j$，$Q''_{i,j}$ 都是前 $j$ 个系数的\Term{polynomial}{多项式}。若 $j<n$，它们就为零，因而 $\gamma^i(x)$ 都属于 $(\operatorname{Ch}A)_n$，于是
\SourcePage{362}
$(\operatorname{Ch}A)_n$ 是\Term{lambda-ideal}{$\lambda$-理想}。
\end{proof}
\begin{proof}[推论 6.6.2 的证明]
设 $x=(n,1+a_1+\cdots+a_n+\cdots)$ 是 $\operatorname{Ch}A$ 中\Term{augmentation}{增广}为 $n$ 的元素；模 $(\operatorname{Ch}A)_{n+1}$，它与 $(n,1+\cdots+a_n)$ 同余，因此 $\lambda_{-1}(x)$ 与
\[
\lambda_{-1}(n,1+\cdots+a_n)=(0,1-(n-1)!a_n+\cdots)
\]
同余，故
\[
\lambda_{-1}(x)=(0,1-(n-1)!a_n+\cdots).
\]
\end{proof}

\begin{SGAresult}{推论}{6.6.4}
设 $\varepsilon$ 是 $\operatorname{Ch}A$ 的\Term{augmentation}{增广}\Term{homomorphism}{同态}；则对任意 $x\in\operatorname{Ch}A$，有
\[
\gamma^n(x-\varepsilon(x))
 =(0,1+(-1)^{n-1}(n-1)!a_n+\cdots).
\]
\end{SGAresult}
\begin{proof}
由定义，
\[
\gamma^n(x-\varepsilon(x))
 =(-1)^n\sum_{i=0}^n(-1)^i\lambda^i(x+n-\varepsilon(x)).
\]
由于 $x+n-\varepsilon(x)$ 的\Term{augmentation}{增广}为 $n$，可以应用推论 6.6.2，所需结论随即成立。
\end{proof}

\begin{SGAresult}{推论}{6.6.5}
沿用 3.10 的记号，有
$\operatorname{Fil}_n\operatorname{Ch}A\subset(\operatorname{Ch}A)_n$。
\end{SGAresult}
\begin{proof}
由上一推论，$\gamma^n(x-\varepsilon(x))\in(\operatorname{Ch}A)_n$。这两个\Term{filtration}{滤过}都是\Term{ring-filtration}{环滤过}，而当 $n$ 和 $x$ 变化时，$\gamma^n(x-\varepsilon(x))$ 生成\Term{lambda-ring}{$\lambda$-环}的\Term{filtration}{滤过}，由此便得到本推论。
\end{proof}

\begin{SGAresult}{推论}{6.6.7}
设 $A$ 是\Term{characteristic}{特征}为零且\Term{degree}{次数}\Term{bounded}{有界}的\Term{graded-algebra}{分次 $K$-代数}；则 $\operatorname{Ch}A$ 上定义的两个\Term{filtration}{滤过}相同。
\end{SGAresult}
\begin{proof}
假设 $\operatorname{Fil}^pA=0$，并设 $x=(0,1+a_n+\cdots+a_p)$ 是 $(\operatorname{Ch}A)_n$ 中的元素。置
\[
x'=\left(0,1+\frac{(-1)^{n-1}}{(n-1)!}a_n+a_{n+1}+\cdots+a_p\right).
\]
则 $\gamma^n(x')=(0,1+a_n+\cdots+a'_p)$，因为 $x'$ 的\Term{augmentation}{增广}为零。因而
\SourcePage{363}
$x-\gamma^n(x')\in(\operatorname{Ch}A)_{n+1}$。将这一操作迭代 $p-n$ 次，就到达为零的 $(\operatorname{Ch}A)_{p+1}$，因而把 $x$ 写成了 $\operatorname{Fil}^n\operatorname{Ch}A$ 中的元素。结合（6.6.5），这两个\Term{filtration}{滤过}便相同。
\end{proof}

\textbf{6.7。} 现在假设 $A$ 是向 $K$ \Term{augmentation}{增广}的\Term{lambda-ring}{$\lambda$-环}，\Term{augmentation}{增广}\Term{homomorphism}{同态}记作 $\varepsilon$。于是 $A$ 带有\Term{augmented-lambda-ring}{增广 $\lambda$-环}的\Term{canonical}{典范}\Term{filtration}{滤过}，从而可以给 $\operatorname{Gr}A$ 关联其\Term{chern-ring}{Chern 环}。

若 $x\in A$，则由\Term{augmented-lambda-ring}{增广 $\lambda$-环}\Term{filtration}{滤过}的定义，
$\gamma_t(x-\varepsilon(x))\in\widehat G(A)$。将 $\gamma_t(x-\varepsilon(x))$ 在
$\widehat G(A)/\widehat G'(A)=1+\widehat{\operatorname{Gr}A}^{+}$ 中的像记作 $c(x)$，称为 $x$ 的\Term{total-chern-class}{总 Chern 类}；以 $c^i(x)$ 表示 $c(x)$ 的第 $i$ 个系数，称为 $x$ 的第 $i$ \Term{chern-class}{Chern 类}。于是
\[
c^i(x)=\operatorname{cl}_{\operatorname{Gr}^iA}
 \bigl(\gamma^i(x-\varepsilon(x))\bigr);
\qquad c(x)=1+c^1(x)+\cdots.
\tag{6.7.1}
\]
再置
\[
\widetilde c(x)=(\varepsilon(x),c(x)),
\tag{6.7.2}
\]
并称 $\widetilde c(x)$ 为 $x$ 的\Term{completed-chern-class}{补全 Chern 类}（classe de Chern complétée）。

\begin{SGAresult}{命题}{6.8}
\Term{completed-chern-class}{补全 Chern 类}
$\widetilde c:A\longrightarrow\operatorname{Ch}(\operatorname{Gr}A)$ 是\Term{augmented-algebra}{增广 $K$-代数}的 $K$-\Term{lambda-homomorphism}{$\lambda$-同态}。
\end{SGAresult}
\begin{proof}
显然，这样定义的\Term{homomorphism}{同态}具有加法性；为验证\Term{multiplication}{乘法}性，写成
\[
\begin{aligned}
\gamma_t(xy-\varepsilon(xy))
={}&\gamma_t\bigl(\varepsilon(y)(x-\varepsilon(x))\bigr)
 \gamma_t\bigl(\varepsilon(x)(y-\varepsilon(y))\bigr)\\
&\hspace{1em}\cdot\gamma_t\bigl((x-\varepsilon(x))(y-\varepsilon(y))\bigr)\\
={}&\gamma_t(x-\varepsilon(x))^{\varepsilon(y)}
 \gamma_t(y-\varepsilon(y))^{\varepsilon(x)}\\
&\hspace{1em}\cdot\bigl(\gamma_t(x-\varepsilon(x))\nabla
 \gamma_t(y-\varepsilon(y))\bigr).
\end{aligned}
\]
\SourcePage{364}
取商后，由（6.1.2）便得到结论。

若 $a\in K$，则 $\widetilde c(a)=(a,1)$，所以 $\widetilde c$ 是 $K$-\Term{algebra}{代数}\Term{homomorphism}{同态}；此外，它与\Term{augmentation}{增广}相容。

最后还需证明它是\Term{lambda-homomorphism}{$\lambda$-同态}。这就要证明下图交换：
\[
\begin{tikzcd}
A \arrow[r,"\gamma_t"] \arrow[d,"\widetilde c"'] &
 \widehat G_g(A) \arrow[d,"\widehat G_g(\widetilde c)"]\\
\operatorname{Ch}A \arrow[r,"\gamma_t"'] &
 \widehat G_g(\operatorname{Ch}A).
\end{tikzcd}
\]
由加法性，只需对 $K$ 中的元素以及 $A$ 中\Term{augmentation}{增广}为零的元素证明：前者显然；对于后者，根据 $c$ 的定义及（6.1.4），所需\Term{commutativity}{交换性}由图（3.7.2）的\Term{commutativity}{交换性}得到。
\end{proof}

\begin{SGAresult}{命题}{6.9}
设 $A$ 是\Term{augmented-k-lambda-algebra}{增广 $K$-$\lambda$-代数}，$z\in\operatorname{Fil}^mA$；则
\[
\gamma^m(z)\equiv(-1)^{m-1}(m-1)!z
 \pmod{\operatorname{Fil}^{m+1}A},
\]
即
\[
c^m(z)=\operatorname{cl}_{\operatorname{Gr}^mA}
 \bigl((-1)^{m-1}(m-1)!z\bigr).
\]
\end{SGAresult}
\begin{proof}
需要证明这两个元素在 $\operatorname{Gr}^mA$ 中的像相同，也就是说，$z$ 的第 $m$ \Term{chern-class}{Chern 类}是 $(-1)^{m-1}(m-1)!z$ 的像。

可以将 $z$ 写成若干项
$\sum a\gamma^{i_1}(x_1)\cdots\gamma^{i_k}(x_k)$ 之和，其中 $a\in A$，每个 $i$ 都有 $\varepsilon(x_i)=0$，且 $i_1+\cdots+i_k\geqslant m$。于是
\[
\widetilde c(z)=\sum\widetilde c(a)\cdot
 \widetilde c(\gamma^{i_1}(x_1))\cdots
 \widetilde c(\gamma^{i_k}(x_k)).
\]
由于 $\widetilde c$ 是\Term{lambda-homomorphism}{$\lambda$-同态}，可以写成
\[
\widetilde c(\gamma^{i_j}(x_j))
 =\gamma^{i_j}(\widetilde c(x_j)).
\]
\SourcePage{365}
结合 $c$ 的定义及（6.6.4），得到
\[
\widetilde c(\gamma^{i_j}(x_j))
 =\bigl(0,1+(-1)^{i_j-1}(i_j-1)!\,\gamma^{i_j}(x_j)+\cdots\bigr).
\]
由（6.5.3）及 $\widetilde c(a)=(a,1)$，最终得到
\[
\begin{aligned}
\widetilde c(z)
 &=\sum\left(0,
 \left(1+(-1)^{i_1+\cdots+i_k-1}(i_1+\cdots+i_k-1)!
 \,\gamma^{i_1}(x_1)\cdots\gamma^{i_k}(x_k)+\cdots\right)^a\right)\\
 &=\left(0,1+
 \sum_{\substack{(x_1,\ldots,x_k)\\i_1+\cdots+i_k=m}}
 (-1)^{m-1}(m-1)!a\,
 \gamma^{i_1}(x_1)\cdots\gamma^{i_k}(x_k)+\cdots\right)\\
 &=\bigl(0,1+(-1)^{m-1}(m-1)!\bar z+\cdots\bigr),
\end{aligned}
\]
这就得到结论。
\end{proof}

\begin{SGAresult}{命题}{6.10}
设 $K$ 是\Term{binomial-ring}{二项式环}，$R$ 是\Term{augmented-k-lambda-algebra}{增广 $K$-$\lambda$-代数}，$N$ 是 $R$ 中满足 $\lambda^k(N)=0$（$k>d$）的元素，且 $x\in\operatorname{Fil}^iR$。则
\[
\gamma^{d+i}(N,x)-(-1)^{d+i-1}(d+i-1)!x
 \in\operatorname{Fil}^{i+1}R.
\]
\end{SGAresult}
\begin{proof}
可以找到元素 $x_{\alpha,j}\in R$ 及 $a_\alpha\in K$，使得
\[
x=\sum_\alpha a_\alpha
 \gamma^{i_1}(x_{\alpha,1})\cdots\gamma^{i_k}(x_{\alpha,k}),
\]
其中 $i_1+\cdots+i_k\geqslant i$，且对任意 $\alpha,j$，都有 $\varepsilon(x_{\alpha,j})=0$。设 $R_0$ 是由\Term{generator}{生成元} $N'_0,X_{\alpha,j}$ 生成的 $K$-\Term{lambda-algebra}{$\lambda$-代数}，满足关系
\[
\lambda^k(N'_0)=0\qquad(k>d).
\]
由 $\varepsilon(N'_0)=\varepsilon(X_{\alpha,j})=0$ 定义\Term{augmentation}{增广}\Term{homomorphism}{同态} $R_0\longrightarrow K$。根据（4.9.1），$R_0$ 是以 $K$ 为系数环、以 $\gamma^n(N'_0)$（$1\leqslant n\leqslant d$）及 $\gamma^p(X_{\alpha,j})$（$p\geqslant1$）为\Term{indeterminate}{不定元}的\Term{polynomial-ring}{多项式环}；若赋予 $\gamma^k(N'_0)$ 和 $\gamma^k(X_{\alpha,j})$ 权 $k$，则 $R_0$ 的 $\lambda$-\Term{filtration}{滤过}与其按权的\Term{filtration}{滤过}相同（4.12 和 4.15）。

置
\[
X_0=\sum_\alpha a_\alpha
 \gamma^{i_1}(X_{\alpha,1})\cdots\gamma^{i_k}(X_{\alpha,k}),
\]
\SourcePage{366}
先对 $X_0$ 和 $N'_0=N_0+d$ 证明 6.10（它们确实满足 6.10 的假设）。考虑到上面关于 $R_0$ 的\Term{filtration}{滤过}的说明，只需证明
\[
\gamma^{d+i}(N_0,X_0)\gamma^d(N'_0)
 -(-1)^{d+i-1}(d+i-1)!X_0\gamma^d(N'_0)
 \in\operatorname{Fil}^{d+i+1}R_0.
\]
而 $\lambda_{-1}(N_0)=(-1)^d\gamma^d(N'_0)$，因此此式可写为
\[
\begin{aligned}
{}&(-1)^d\gamma^{d+i}\bigl(X_0(-1)^d\gamma^d(N'_0)\bigr)\\
&\qquad-(-1)^{d+i-1}(d+i-1)!X_0\gamma^d(N'_0)
 \in\operatorname{Fil}^{d+i+1}R_0.
\end{aligned}
\]
置 $z=(-1)^dX_0\gamma^d(N'_0)$；则 $z\in\operatorname{Fil}^{d+i}R_0$，从而归结为证明
\[
\gamma^{d+i}(z)-(-1)^{d+i-1}(d+i-1)!z
 \in\operatorname{Fil}^{d+i+1}R_0.
\]
这由 6.9 得到。

为由此特例\Term{pushout}{推出} 6.10，考虑将 $N'_0$ 映到 $N-d$、将每个 $X_{\alpha,j}$ 映到 $x_{\alpha,j}$ 的 $K$-\Term{lambda-homomorphism}{$\lambda$-同态} $R_0\longrightarrow R$（4.7）。这个\Term{homomorphism}{同态}与 $\lambda$-\Term{filtration}{滤过}相容，并将 $N_0$ 映到 $N$、将 $X_0$ 映到 $x$，因而得到结论。
\end{proof}

\begin{SGAresult}{命题}{6.11}
假设 $K$ 的\Term{characteristic}{特征}为零，设 $n\geqslant0$ 是整数。则\Term{functor}{函子} $\operatorname{Ch}$ 给出如下\Term{equivalence-of-categories}{范畴等价}：从满足 $A_0=K$、$A_k=0$（$k>n$）的\Term{graded-algebra}{分次 $K$-代数} $A$ 的\Term{category}{范畴}，到满足 $\operatorname{Fil}^{n+1}\Lambda=0$ 的\Term{augmented-k-lambda-algebra}{增广 $K$-$\lambda$-代数} $\Lambda$ 的\Term{category}{范畴}；一个\Term{quasi-inverse-functor}{拟逆函子}由 $\Lambda\longmapsto\operatorname{Gr}\Lambda$ 给出。
\end{SGAresult}
\begin{proof}
设 $A$ 是满足 $A_k=0$（$k>n$）的\Term{graded-algebra}{分次 $K$-代数}，带有其\Term{grading}{分次}所关联的\Term{filtration}{滤过}。于是 $(\operatorname{Ch}A)_{n+1}=0$，故 $\operatorname{Fil}^{n+1}\operatorname{Ch}A=0$（6.6.5）。反过来，若 $\Lambda$ 是满足 $\operatorname{Fil}^{n+1}\Lambda=0$ 的\Term{augmented-k-lambda-algebra}{增广 $K$-$\lambda$-代数}，则 $\operatorname{Gr}^0\Lambda=K$，$\operatorname{Gr}^k\Lambda=0$（$k>n$）。

下面定义\Term{functor}{函子}的\Term{isomorphism}{同构}
\[
A\longrightarrow\operatorname{Gr}(\operatorname{Ch}A),
\qquad
\Lambda\longrightarrow\operatorname{Ch}(\operatorname{Gr}\Lambda).
\]
\SourcePage{367}
由于 $A$ 的\Term{grading}{分次}\Term{finite}{有限}，$\operatorname{Ch}A$ 的\Term{lambda-ring}{$\lambda$-环}\Term{filtration}{滤过}和\Term{degree-filtration}{次数滤过}的\Term{associated-graded-ring}{相伴分次环}相同（6.6.7）。于是定义
$u:A\longrightarrow\operatorname{Gr}_{\deg}(\operatorname{Ch}A)$：在 $A_0$ 上取\Term{identity-map}{恒等映射}，并令
\[
\forall i\geqslant1,\ \forall a_i\in A_i,\qquad
u(a_i)=\operatorname{cl}_{\operatorname{Gr}^i}
 \bigl(0,1+(-1)^{i-1}(i-1)!a_i\bigr),
\]
再按加法性扩充。根据关系式（6.5.3），它是\Term{graded-ring}{分次环}的\Term{homomorphism}{同态}，并且容易看出它与 $K$ 的运算相容。逆\Term{morphism}{态射} $u^{-1}$ 可以如下定义：对任意 $i\geqslant1$、$\bar\varphi\in\operatorname{Gr}^i(\operatorname{Ch}A)$，$\bar\varphi$ 的第 $i$ 个系数 $a_i$ 不依赖于将 $\bar\varphi$ \Term{lifting}{提升}到 $(\operatorname{Ch}A)_i$ 的方式，置
\[
u^{-1}(\bar\varphi)=\frac{(-1)^{i-1}}{(i-1)!}a_i.
\]
因此\Term{homomorphism}{同态} $u$ 是\Term{isomorphism}{同构}。

\Term{isomorphism}{同构} $\Lambda\longrightarrow\operatorname{Ch}(\operatorname{Gr}\Lambda)$ 由\Term{completed-chern-class}{补全 Chern 类} $\widetilde c$ 定义。由于\Term{filtration}{滤过}是\Term{discrete-filtration}{离散的}，要证明它确实定义了\Term{isomorphism}{同构}，只需证明取\Term{associated-graded-ring}{相伴分次环}所得的\Term{homomorphism}{同态}是\Term{isomorphism}{同构}。得到
\[
\begin{aligned}
\operatorname{Gr}\Lambda
 &\xrightarrow{\operatorname{Gr}(\widetilde c)}
 \operatorname{Gr}_{\lambda}(\operatorname{Ch}(\operatorname{Gr}\Lambda))\\
 &\simeq\operatorname{Gr}_{\deg}(\operatorname{Ch}(\operatorname{Gr}\Lambda))
 \xrightarrow{u^{-1}}\operatorname{Gr}\Lambda.
\end{aligned}
\]
只需证明其复合是\Term{identity-map}{恒等映射}。由于这些\Term{homomorphism}{同态}是\Term{graded-algebra}{分次 $K$-代数}的\Term{homomorphism}{同态}，根据 $\Lambda$ 的\Term{lambda-ring}{$\lambda$-环}\Term{filtration}{滤过}的定义，只需对满足 $\varepsilon(x)=0$ 的 $\overline{\gamma^i(x)}$ 证明这一点。得到
\[
\operatorname{Gr}(\widetilde c)(\overline{\gamma^i(x)})
 =\operatorname{cl}_{\operatorname{Gr}^i}
 \bigl(\gamma^i(\widetilde c(x))\bigr).
\]
由（6.6.4），
\[
\gamma^i(\widetilde c(x))
 =(0,1+(-1)^{i-1}(i-1)!c^i(x)+\cdots),
\]
所以
\[
\operatorname{cl}_{\operatorname{Gr}^i}
 \bigl(\gamma^i(\widetilde c(x))\bigr)=u(c^i(x)),
\]
\SourcePage{368}
这说明 $u^{-1}\circ\operatorname{Gr}(\widetilde c)=\operatorname{Id}$。
\end{proof}

\begin{SGAresult}{推论}{6.12}
若 $K$ 的\Term{characteristic}{特征}为零，且 $\Lambda$ 是带有\Term{discrete-filtration}{离散滤过}的 $K$-\Term{lambda-algebra}{$\lambda$-代数}，则\Term{completed-chern-class}{补全 Chern 类}
\[
\widetilde c:\Lambda\longrightarrow\operatorname{Ch}(\operatorname{Gr}\Lambda)
\]
是\Term{isomorphism}{同构}。
\end{SGAresult}
\begin{proof}
证明已包含在 6.11 的证明中。
\end{proof}

\begin{SGAresult}{定义}{6.13}
设 $A$ 是\Term{graded-algebra}{分次 $K$-代数}，$\Lambda$ 是\Term{augmented-k-lambda-algebra}{增广 $K$-$\lambda$-代数}。称每个与\Term{augmentation}{增广}相容的\Term{lambda-homomorphism}{$\lambda$-同态}
$\Lambda\longrightarrow\operatorname{Ch}(A)$ 为 $\Lambda$ 的取值于 $A$ 的\Term{chern-theory}{Chern 理论}。
\end{SGAresult}
给每个\Term{graded-algebra}{分次 $K$-代数} $A$ 关联 $\Lambda$ 的取值于 $A$ 的\Term{chern-theory}{Chern 理论}所成的\Term{set}{集合}，即
$\operatorname{Hom}_{\lambda}(\Lambda,\operatorname{Ch}(A))$，便在\Term{graded-algebra}{分次 $K$-代数}的\Term{category}{范畴}上定义了一个取值于\Term{set}{集合}\Term{category}{范畴}的\Term{functor}{函子}。

\begin{SGAresult}{命题}{6.14}
6.13 中定义的\Term{functor}{函子}是\Term{representable-functor}{可表的}。若 $K$ 的\Term{characteristic}{特征}为零，且 $\Lambda$ 的\Term{filtration}{滤过}是\Term{discrete-filtration}{离散的}，则它由 $\Lambda$ 的\Term{associated-graded-ring}{相伴分次环}表示。
\end{SGAresult}
\begin{proof}
先证明第一点。考虑由\Term{generator}{生成元} $c^i(x)$（$i\geqslant1$、$x\in\operatorname{Fil}^1\Lambda$）生成的\Term{graded-algebra}{分次 $K$-代数}，其中 $c^i(x)$ 的\Term{degree}{次数}为 $i$；按以下关系取商：
\[
\begin{aligned}
\textup{i)}\quad
c^k(x+y)&=\sum_{i+j=k}c^i(x)c^j(y),
 &&k\geqslant1,\ x,y\in\operatorname{Fil}^1\Lambda;\\
\textup{ii)}\quad
c^k(xy)&=P''_k(c^i(x),c^j(y))\quad\text{（参见 6.1）},
 &&k\geqslant1,\ x,y\in\operatorname{Fil}^1\Lambda;\\
\textup{iii)}\quad
c^k(ax)&=R_{k,a}(c^i(x))\quad\text{（参见 6.1）},
 &&k\geqslant1,\ a\in K,\ x\in\operatorname{Fil}^1\Lambda;\\
\textup{iv)}\quad
c^k(\gamma^j(x))&=Q''_{j,k}(c^i(x))\quad\text{（参见 6.1）},
 &&j,k\geqslant1,\ x\in\operatorname{Fil}^1\Lambda.
\end{aligned}
\]
\SourcePage{369}
由于这些关系是\Term{isobaric-polynomial}{等权的}，所得商是一个\Term{graded-algebra}{分次 $K$-代数} $\Omega$。定义\Term{morphism}{态射}
$\widetilde c_0:\Lambda\longrightarrow\operatorname{Ch}(\Omega)$ 为
\[
x\longmapsto\bigl(\varepsilon(x),1+c^1(x-\varepsilon(x))+\cdots\bigr).
\]
所加的关系说明，这是 $\Lambda$ 的取值于 $\Omega$ 的\Term{chern-theory}{Chern 理论}。此外，若有一个取值于\Term{graded-algebra}{分次 $K$-代数} $A$ 的\Term{chern-theory}{Chern 理论} $f$，则存在唯一的\Term{homomorphism}{同态} $\Omega\longrightarrow A$，使 $f$ 通过 $\operatorname{Ch}(\Omega)$ \Term{resolution}{分解}；其定义是：对每个 $i\geqslant1$ 及 $x\in\operatorname{Fil}^1\Lambda$，将 $c^i(x)$ 映到 $f(x)$ 的第 $i$ 个系数（因为 $f$ 是\Term{lambda-homomorphism}{$\lambda$-同态}，这确实是良定义的）。

现在假设 $K$ 的\Term{characteristic}{特征}为零，且 $\Lambda$ 的\Term{filtration}{滤过}是\Term{discrete-filtration}{离散的}；取 $n$ 使 $\operatorname{Fil}^n\Lambda=0$。则\Term{graded-algebra}{分次代数} $\Omega$ 在\Term{degree}{次数} $\geqslant n$ 处为零。为此，需要证明：若 $x_1,\ldots,x_p\in\operatorname{Fil}^1\Lambda$ 且 $k_1+\cdots+k_p\geqslant n$，则
\[
c^{k_1}(x_1)\cdots c^{k_p}(x_p)=0.
\]
而由假设，
\[
\gamma^{k_1}(x_1)\cdots\gamma^{k_p}(x_p)=0.
\]
应用\Term{lambda-homomorphism}{$\lambda$-同态} $\widetilde c_0$，得到
\[
\gamma^{k_1}(\widetilde c_0(x_1))\cdots
\gamma^{k_p}(\widetilde c_0(x_p))=(0,1).
\]
根据（6.6.4）和（6.5.3），这\Term{pushout}{推出}
\[
(k_1+\cdots+k_p)!c^{k_1}(x_1)\cdots c^{k_p}(x_p)=0.
\]
由于 $K$ 的\Term{characteristic}{特征}为零，便得到所需结论。

将 $\Omega$ 的\Term{universal-property}{泛性质}应用于\Term{canonical}{典范}\Term{homomorphism}{同态} $\widetilde c$，得到\Term{commutative-diagram}{交换图}
\SourcePage{370}
\[
\begin{tikzcd}
\Lambda \arrow[r,"\widetilde c"] \arrow[d,"\widetilde c_0"'] &
 \operatorname{Ch}(\operatorname{Gr}\Lambda)\\
\operatorname{Ch}(\Omega) \arrow[ur] &
\end{tikzcd}
\]
由命题 6.11，\Term{homomorphism}{同态}
$\Omega\longrightarrow\operatorname{Gr}\Lambda$ 是\Term{isomorphism}{同构}，当且仅当
$\operatorname{Ch}(\Omega)\longrightarrow\operatorname{Ch}(\operatorname{Gr}\Lambda)$ 是\Term{isomorphism}{同构}。前一个显然是\Term{surjectivity}{满射}，所以后一个也是\Term{surjectivity}{满射}，因而只需证明
$\operatorname{Ch}(\Omega)\longrightarrow\operatorname{Ch}(\operatorname{Gr}\Lambda)$ 为\Term{injectivity}{单射}。由 $\operatorname{Ch}$ 的\Term{functoriality}{函子性}，得到图
\[
\begin{tikzcd}
\Lambda \arrow[r,"\widetilde c_0"] \arrow[d,"\widetilde c"'] &
 \operatorname{Ch}(\Omega) \arrow[d,"\widetilde c"]\\
\operatorname{Ch}(\operatorname{Gr}\Lambda)
 \arrow[r,"\operatorname{Ch}(\widetilde c_0)"'] &
 \operatorname{Ch}(\operatorname{Gr}_{\lambda}(\operatorname{Ch}(\Omega))).
\end{tikzcd}
\]
由于 $\Omega$ 的\Term{degree}{次数}\Term{bounded}{有界}，右侧竖\Term{arrow}{箭头}是\Term{isomorphism}{同构}（6.12）；再利用 6.11，便得到一个\Term{morphism}{态射}
$\operatorname{Gr}\Lambda\longrightarrow\Omega$。根据 $\Omega$ 的\Term{universal-property}{泛性质}，两者的复合是\Term{identity-map}{恒等映射}。
\end{proof}

\SGAsection{}{附录：Adams 运算 $\psi^k$}
这里采用 J.-P. Serre 给出的一种叙述方式。

\textbf{7.1。} 设 $K$ 是\Term{pre-lambda-ring}{预 $\lambda$-环}。对每个 $x\in K$，以
$\frac{d}{dt}(\lambda_t(x))$ 表示 $\lambda_t(x)$ 关于 $t$ 的导数。定义 $K$ 中的元素 $\psi^k(x)$，使
\[
\frac{d}{dt}(\lambda_t(x))/\lambda_t(x)
 =\sum_{k=1}^{\infty}(-1)^{k-1}\psi^k(x)t^{k-1}.
\tag{7.1.1}
\]
映射 $\psi^k:K\longrightarrow K$ 是\Term{additive}{加性}的。事实上，
\[
\begin{aligned}
\frac{d}{dt}(\lambda_t(x+y))/\lambda_t(x+y)
 &=\frac{d}{dt}(\lambda_t(x)\lambda_t(y))
   /\bigl(\lambda_t(x)\lambda_t(y)\bigr)\\
 &=\frac{d}{dt}(\lambda_t(x))/\lambda_t(x)
  +\frac{d}{dt}(\lambda_t(y))/\lambda_t(y),
\end{aligned}
\]
\SourcePage{371}
因而对每个 $k$，都有
\[
\psi^k(x+y)=\psi^k(x)+\psi^k(y).
\]
由（7.1.1），$\psi^k$ 是关于 $\lambda^i$（$i=1,\ldots,k$）的权 $k$ 的\Term{isobaric-polynomial}{等权多项式}。例如，$\psi^1$ 和 $\psi^2$ 的值为
\[
\begin{aligned}
\psi^1(x)&=\lambda^1(x)=x,\\
\psi^2(x)&=x^2-2\lambda^2(x).
\end{aligned}
\]

\textbf{7.2。} 设 $K$ 是\Term{characteristic}{特征}为 $0$ 的环；假设给定\Term{additive-map}{加性映射} $\psi^k:K\longrightarrow K$，满足 $\psi^1=\operatorname{Id}_K$。置
\[
\lambda_t(x)=\exp\left(\sum_{k=1}^{\infty}
 (-1)^{k-1}\psi^k(x)t^k/k\right),
\tag{7.2.1}
\]
便在 $K$ 上定义了一个\Term{pre-lambda-structure}{预 $\lambda$-结构}。\Term{multiplication}{乘法}性由各 $\psi^k$ 的加法性得到。因此，在 $K$ 上给定一个\Term{pre-lambda-structure}{预 $\lambda$-结构}，\Term{equivalence}{等价}于给定相应的\Term{adams-operation}{Adams 运算} $\psi^k$。

\textbf{7.3。} 由公式（7.1.1）显然可知，$\psi^k$ 与\Term{lambda-homomorphism}{$\lambda$-同态}交换。若 $K,K'$ 是\Term{characteristic}{特征}为零的\Term{pre-lambda-ring}{预 $\lambda$-环}，则一个\Term{homomorphism}{同态} $f:K\longrightarrow K'$ 是\Term{lambda-homomorphism}{$\lambda$-同态}，当且仅当它与所有 $\psi^k$ 交换。这由 7.2 得到，因为各 $\lambda^i$ \Term{representable-functor}{可表示}为关于 $\psi^k$ 的\Term{polynomial}{多项式}。

\textbf{例 7.4。} i）在 $K=\mathbb Z$ 上取其\Term{lambda-ring}{$\lambda$-环}结构，则
\[
\frac{d}{dt}(\lambda_t(n))/\lambda_t(n)
 =\frac{n(1+t)^{n-1}}{(1+t)^n}
 =n\sum_{k=1}^{\infty}(-1)^{k-1}t^{k-1}.
\]
因此，所有 $\psi^k$ 都等于 $\mathbb Z$ 的\Term{identity-map}{恒等映射}。

\textup{ii)} 考虑 $\mathbb Z$-\Term{algebra}{代数}的\Term{category}{范畴} $\mathscr C_0$（1.8）。由 7.3，
\SourcePage{372}
这些\Term{adams-operation}{Adams 运算} $\psi^k$ 定义了取值于\Term{lambda-ring}{$\lambda$-环}的\Term{functor}{函子} $\widehat G$ 的\Term{additive}{加性}\Term{endomorphism}{自同态}；由 1.10，它们由在 $1+Xt$ 上的值确定。有
\[
\begin{aligned}
\frac{d}{dt}(\lambda_u(1+Xt))/\lambda_u(1+Xt)
 &=\frac{1+Xt}{1+\{1+Xt\}u}\\
 &=\sum_{k=1}^{\infty}(-1)^{k-1}
 \{1+Xt\}^{\circ k}u^{k-1}.
\end{aligned}
\]
所以
\[
\psi^k(1+Xt)=1+X^kt.
\tag{7.4.1}
\]

\begin{SGAresult}{命题}{7.5}
\textup{i)} 若 $K$ 是\Term{lambda-ring}{$\lambda$-环}，则\Term{adams-operation}{Adams 运算} $\psi^k$ 满足恒等式
\[
\psi^k(1)=1;\qquad
\psi^k(xy)=\psi^k(x)\psi^k(y);\qquad
\psi^k\circ\psi^{k'}=\psi^{kk'}.
\tag{7.5.1}
\]
\textup{ii)} 反过来，若各 $\psi^k$ 满足这些恒等式，且 $K$ 的\Term{characteristic}{特征}为零，则 $K$ 是\Term{lambda-ring}{$\lambda$-环}。
\end{SGAresult}
\begin{proof}
i）若 $K$ 是\Term{lambda-ring}{$\lambda$-环}，则 $\lambda_t(1)=1+t$，从而
\[
\frac{d}{dt}(\lambda_t(1))/\lambda_t(1)
 =\frac{1}{1+t}
 =\sum_{k=1}^{\infty}(-1)^{k-1}t^{k-1},
\]
因而 $\psi^k(1)=1$。

为证明关系式（7.5.1）中的第二式，考虑图
\[
\begin{tikzcd}[column sep=large]
K\times K \arrow[r,"\lambda_t\times\lambda_t"] \arrow[d] &
 \widehat G(K)\times\widehat G(K)
 \arrow[r,"\Phi^k\times\Phi^k"] \arrow[d,"\circ"] &
 K\times K \arrow[d]\\
K \arrow[r,"\lambda_t"'] &
 \widehat G(K) \arrow[r,"\Phi^k"'] & K,
\end{tikzcd}
\tag{7.5.2}
\]
其中竖\Term{arrow}{箭头}由\Term{multiplication}{乘法}定义，而映射 $\Phi^k$ 如下定义。“\Term{logarithmic-derivative}{对数导数}”映射是
\SourcePage{373}
从 $\widehat G(K)$ 到 $K[[t]]$ 的群\Term{homomorphism}{同态}；将一个级数映到其第 $k-1$ 个系数乘以 $(-1)^{k-1}$，是从 $K[[t]]$ 到 $K$ 的\Term{additive}{加性}\Term{homomorphism}{同态}；$\Phi^k$ 定义为两者的复合。根据定义（7.1.1），显然有 $\Phi^k\circ\lambda_t=\psi^k$。

因此，公式（7.5.1）的第二式正表示上图的\Term{commutativity}{交换性}。若假定 $K$ 是\Term{lambda-ring}{$\lambda$-环}，则左侧方块交换，只需验证右侧方块交换。这是 $\mathscr C_0$ 上以 $\widehat G\times\widehat G$ 为源的\Term{functor}{函子}\Term{homomorphism}{同态}的图；根据 1.10，只需对 $(1+Xt,1+Yt)$ 验证。有
\[
\Phi^k(1+Xt)=X^k
 \quad\text{（$X/(1+Xt)$ 的第 $k-1$ 个系数）};
\qquad \Phi^k(1+Yt)=Y^k;
\]
\[
\begin{aligned}
\Phi^k\bigl((1+Xt)\circ(1+Yt)\bigr)
 &=\Phi^k(1+XYt)=(XY)^k\\
 &=\Phi^k(1+Xt)\Phi^k(1+Yt),
\end{aligned}
\]
这就得到结论。

为证明最后一个关系式，先指出下图交换：
\[
\begin{tikzcd}
\widehat G^t(K) \arrow[r,"\lambda_u"] \arrow[d,"\Phi^{kk'}"'] &
 \widehat G^u\circ\widehat G^t(K) \arrow[d,"\Phi_u^{k'}"]\\
K \arrow[r,"\Phi^k"'] & \widehat G^t(K).
\end{tikzcd}
\tag{7.5.3}
\]
如上，只需对 $1+Xt$ 验证，而
\[
\begin{aligned}
\Phi_u^{k'}(\lambda_u(1+Xt))
 &=\Phi_u^{k'}(1+\{1+Xt\}u)
 =(1+Xt)^{\circ k'}=1+X^{k'}t;\\
\Phi^k(1+X^{k'}t)&=X^{kk'}=\Phi^{kk'}(1+Xt).
\end{aligned}
\]
所以
\[
\psi^{kk'}
 =\Phi^k\circ\Phi_u^{k'}\circ\lambda_u\circ\lambda_t
 =\Phi^k\circ\psi^{k'}\circ\lambda_t.
\]
另一方面，
\[
\psi^k\circ\psi^{k'}
 =\Phi^k\circ\lambda_t\circ\Phi^{k'}\circ\lambda_t
 =\Phi^k\circ\lambda_t\circ\psi^{k'}.
\]
\SourcePage{374}
若 $K$ 是\Term{lambda-ring}{$\lambda$-环}，则图
\[
\begin{tikzcd}
K \arrow[r,"\lambda_t"] \arrow[d,"\psi^{k'}"'] &
 \widehat G(K) \arrow[d,"\psi^{k'}"]\\
K \arrow[r,"\lambda_t"'] & \widehat G(K)
\end{tikzcd}
\tag{7.5.4}
\]
由 7.3 是交换的，因为 $\lambda_t$ 是\Term{lambda-homomorphism}{$\lambda$-同态}。所需关系式随即得到。

\textup{ii)} 若 $K$ 的\Term{characteristic}{特征}为零，则“\Term{logarithmic-derivative}{对数导数}”映射是从 $\widehat G(K)$ 到 $K[[t]]$ 的\Term{isomorphism}{同构}，因为可以积分。因此，$\widehat G(K)$ 中两个元素相等，当且仅当对每个 $k$，它们经 $\Phi^k$ 所得的像相同。

要证明 $\lambda_t$ 是环\Term{homomorphism}{同态}，须证明图（7.5.2）的左侧方块交换。根据上一说明，只需对每个 $k$ 证明图（7.5.2）交换，因为右侧方块总是交换的。而这恰好就是关系式
$\psi^k(xy)=\psi^k(x)\psi^k(y)$。

要证明 $\lambda_t$ 是\Term{lambda-homomorphism}{$\lambda$-同态}，由 7.3，只需对每个 $k'$ 证明图（7.5.4）交换。根据上面的说明，这\Term{equivalence}{等价}于对任意 $k,k'$ 证明
\[
\Phi^k\circ\psi^{k'}\circ\lambda_t
 =\Phi^k\circ\lambda_t\circ\psi^{k'},
\]
即关系式 $\psi^{kk'}=\psi^k\circ\psi^{k'}$。

最后，由关系式 $\psi^k(1)=1$ 积分得到 $\lambda_t(1)=1+t$。
\end{proof}

\SGAsection{}{参考文献}
\begin{enumerate}
\item[{[1]}] A. Grothendieck，\emph{Chern 类理论}（La théorie des classes de Chern），Bull. Soc. math. France，86，1958，第 137--154 页。
\item[{[2]}] F. Hirzebruch，\emph{\Term{algebraic-geometry}{代数几何}中的新\Term{topology}{拓扑}方法}（Neue topologische Methoden in der algebraischen Geometrie）。
\end{enumerate}
